\documentclass[12pt]{article}

\usepackage{sbc-template}
\usepackage{graphicx,url}
\usepackage[utf8]{inputenc}
\usepackage[brazilian]{babel}
\usepackage{amsmath}
\usepackage{amssymb}
\usepackage{booktabs}
\usepackage{multirow}
\usepackage{array}
\usepackage{xcolor}
\usepackage{colortbl}
\usepackage{tikz}
\usetikzlibrary{shapes.geometric, arrows.meta, positioning, fit, backgrounds}
\usepackage{enumitem}
\setlist{itemsep=0.5em}
\usepackage{indentfirst}
\usepackage{float}

\sloppy

\title{O Efeito do Escalonamento do Meta-Alvo na Performance de Sistemas de Meta-Aprendizagem para Seleção de Algoritmos}

\author{Lael de Lima Santa Rosa\inst{1}}

\address{Instituto de Computação -- Universidade Federal de Alagoas (UFAL)\\
  Maceió -- AL -- Brasil
  \email{llsr@ic.ufal.br}
}

\begin{document}

\maketitle
\begin{resumo}
A seleção de algoritmos em aprendizado de máquina frequentemente utiliza meta-aprendizagem, onde um meta-modelo aprende a prever o desempenho de algoritmos base a partir das características dos conjuntos de dados (meta-features). No entanto, os escores de desempenho que compõem o meta-alvo variam em escala e distribuição, levantando a questão de se técnicas de escalonamento aplicadas ao meta-alvo afetam o meta-modelo. Inspirados por Amorim et al. (2022), que demonstraram que a escolha do escalonamento importa para classificadores de nível base, este trabalho transpõe a questão para o nível meta. Avaliamos quatro meta-modelos (Ridge, SVR, Random Forest, MLP) combinados com cinco estratégias de escalonamento (Baseline, StandardScaler, MinMaxScaler, RobustScaler, QuantileTransformer) em quatro cenários experimentais: 50 conjuntos de dados controlados, 144 datasets reais do OpenML com 8 meta-features manuais, 200 datasets reais com 8 meta-features manuais, e 200 datasets válidos com 78 meta-features extraídas automaticamente pelo PyMFE. Os resultados mostram que a sensibilidade ao escalonamento é determinada pela arquitetura: o MLP colapsa sem escalonamento centralizado (Spearman cai a 0,03 com meta-features em dados reais), o SVR se beneficia para tarefas de ranqueamento, o Random Forest é invariante e o QuantileTransformer degrada consistentemente a capacidade de ranqueamento. Meta-features mais ricas reduzem substancialmente o erro de predição em todos os meta-modelos. Nenhum escalonador domina universalmente.
\end{resumo}

% ============================================================
\section{Introdução}
% ============================================================

A seleção de algoritmos de aprendizado de máquina (ML) é uma das tarefas mais custosas na construção de sistemas inteligentes. Em vez de testar exaustivamente todos os candidatos para um novo dataset --- inviável em tempo e recursos ---, a \textit{meta-aprendizagem} usa experiências passadas para treinar um meta-modelo que mapeia características dos datasets (\textit{meta-features}) ao desempenho real dos algoritmos base (\textit{meta-alvo}), permitindo recomendar a melhor opção para um novo cenário~\cite{vanschoren2018meta}.

Embora as \textit{meta-features} de entrada normalmente passem por pré-processamento cuidadoso, os escores de desempenho que compõem o \textit{meta-alvo} costumam ser usados em sua forma bruta. Métricas como o F1-macro, porém, variam muito em escala e distribuição, e essa amplitude pode prejudicar meta-modelos sensíveis a essas variações. O trabalho de \cite{amorim2022choice} demonstrou, no nível base, que \textit{``a escolha da técnica de escalonamento importa para a performance de classificação''} --- este trabalho transpõe a mesma questão para o nível meta.

\textbf{Problema:} O escalonamento do meta-alvo em sistemas de meta-aprendizagem para seleção de algoritmos produz diferenças relevantes de performance? Existe uma recomendação geral que funcione para qualquer meta-modelo?

\textbf{Hipóteses:} (H1) Meta-modelos sensíveis ao escalonamento apresentarão diferença estatisticamente significativa entre as técnicas. (H2) Meta-modelos baseados em árvores serão insensíveis. (H3) Não existirá uma técnica universalmente superior.

\textbf{Contribuições principais:}
\begin{itemize}
  \item Primeira investigação sistemática do impacto do escalonamento do \textit{meta-alvo} (e não das meta-features) em sistemas de meta-aprendizagem;
  \item Demonstração de que o MLP, sem escalonamento centralizado, produz recomendações próximas ao acaso com dados reais, resultado revertido pelo StandardScaler ou RobustScaler;
  \item Identificação do QuantileTransformer como prejudicial à capacidade de ranqueamento na maioria dos meta-modelos;
  \item Validação em quatro cenários experimentais: (i) dados controlados; (ii) 144 datasets reais do OpenML; (iii) 200 datasets reais; e (iv) 200 datasets com 78 meta-features extraídas pelo PyMFE;
  \item Demonstração de que meta-features ricas (PyMFE) reduzem substancialmente o erro de predição.
\end{itemize}

% ============================================================
\section{Proposta}
% ============================================================

\subsection{Formulação do Problema}

Seja $\mathcal{D} = \{d_1, \dots, d_N\}$ um conjunto de $N$ datasets de classificação. Para cada $d_i$, extraímos um vetor de meta-features $\mathbf{x}_i \in \mathbb{R}^D$, formando a matriz $X \in \mathbb{R}^{N \times D}$. Avaliamos $M$ algoritmos base $\mathcal{A} = \{a_1, \dots, a_M\}$ em cada dataset, obtendo o vetor de performances $\mathbf{y}_i \in \mathbb{R}^M$ (F1-macro de cada $a_j$ em $d_i$) e a matriz de meta-alvo $Y \in \mathbb{R}^{N \times M}$. O objetivo é treinar uma função $f$ que preveja $\hat{\mathbf{y}}_i \approx \mathbf{y}_i$ a partir de $\mathbf{x}_i$, recomendando o algoritmo com melhor estimativa para um novo dataset.

\subsection{Pipeline de Meta-Aprendizagem com Escalonamento do Alvo}

A proposta central é estudar o efeito de uma transformação $S: \mathbb{R}^{N \times M} \rightarrow \mathbb{R}^{N \times M}$ aplicada sobre $Y$ \textit{antes} do treinamento. A Figura~\ref{fig:pipeline} ilustra o pipeline proposto.

\begin{figure}[ht]
\centering
\begin{tikzpicture}[
  node distance=0.5cm and 0.4cm,
  box/.style={draw, rounded corners, minimum width=2.0cm, minimum height=0.7cm,
              font=\scriptsize, align=center, fill=white},
  databox/.style={box, fill=blue!10},
  procbox/.style={box, fill=orange!15},
  modelbox/.style={box, fill=green!15},
  evalbox/.style={box, fill=red!10},
  arr/.style={-{Stealth[scale=0.7]}, thick},
  phase/.style={draw=gray!50, dashed, rounded corners, inner sep=4pt}
]

% Nodes – Train phase
\node[databox]  (datasets)  {Datasets\\ $\mathcal{D}$};
\node[procbox,  right=of datasets]  (mf)    {Meta-features\\ $X$};
\node[procbox,  right=of mf]        (perf)  {Avaliação\\ base $Y$};
\node[procbox,  right=of perf]      (scale) {Scaler $S$\\ sobre $Y_{train}$};
\node[modelbox, right=of scale]     (train) {Treina\\ $f(X \!\to\! Y_s)$};

% Nodes – Test phase (below)
\node[databox,  below=1.0cm of datasets] (newds)  {Novo\\ dataset};
\node[procbox,  right=of newds]          (newmf)  {Meta-features\\ $\mathbf{x}_{new}$};
\node[modelbox, right=of newmf]          (pred)   {Predição\\ $\hat{\mathbf{y}}_s$};
\node[procbox,  right=of pred]           (inv)    {Inversão\\ $S^{-1}(\hat{\mathbf{y}}_s)$};
\node[evalbox,  right=of inv]            (rec)    {Recomenda\\ $a^*$};

% Arrows – Train
\draw[arr] (datasets) -- (mf);
\draw[arr] (mf)       -- (perf);
\draw[arr] (perf)     -- (scale);
\draw[arr] (scale)    -- (train);

% Arrows – Test
\draw[arr] (newds)  -- (newmf);
\draw[arr] (newmf)  -- (pred);
\draw[arr] (train)  -- (pred);
\draw[arr] (pred)   -- (inv);
\draw[arr] (inv)    -- (rec);

% Phase boxes
\begin{scope}[on background layer]
  \node[phase, fit=(datasets)(mf)(perf)(scale)(train),
        label={[font=\scriptsize\itshape]above:Fase de Treino}] {};
  \node[phase, fit=(newds)(newmf)(pred)(inv)(rec),
        label={[font=\scriptsize\itshape]below:Fase de Teste/Inferência}] {};
\end{scope}
\end{tikzpicture}
\caption{Pipeline de meta-aprendizagem com escalonamento do meta-alvo. O escalonador $S$ é ajustado exclusivamente nos dados de treino, evitando \textit{data leakage}. A inversão $S^{-1}$ garante que as métricas finais sejam avaliadas na escala original do F1-score.}
\label{fig:pipeline}
\end{figure}

\subsection{Técnicas de Escalonamento Avaliadas}

As cinco técnicas, aplicadas coluna a coluna sobre $Y$, são:
\begin{enumerate}
  \item \textbf{Baseline} (sem escalonamento): $\tilde{y} = y$;
  \item \textbf{StandardScaler} (Z-score): $\tilde{y} = (y - \mu) / \sigma$;
  \item \textbf{MinMaxScaler}: $\tilde{y} = (y - y_{\min}) / (y_{\max} - y_{\min})$;
  \item \textbf{RobustScaler}: $\tilde{y} = (y - \text{mediana}) / \text{IQR}$, resistente a outliers;
  \item \textbf{QuantileTransformer}: $\tilde{y} = \Phi^{-1}\!\left( \frac{rank(y)}{n+1} \right)$, transformação não-linear que força distribuição normal.
\end{enumerate}

As técnicas 1--4 são \textbf{monotônicas} (preservam a ordenação relativa), enquanto o QuantileTransformer é \textbf{não-linear} e pode distorcer a hierarquia de ranks entre algoritmos.

\subsection{Métricas de Avaliação}

O meta-modelo é avaliado em três dimensões:
\begin{itemize}
  \item \textbf{MAE}: erro absoluto médio na escala original do F1-score;
  \item \textbf{Spearman} ($\rho_s$): correlação entre a ordenação real e a prevista dos $M$ algoritmos; varia de $-1$ a $+1$;
  \item \textbf{Acc@1}: proporção de datasets em que o algoritmo mais recomendado é de fato o melhor.
\end{itemize}

% ============================================================
\section{Metodologia Experimental}
% ============================================================

\subsection{Conjuntos de Dados}

Foram conduzidos quatro experimentos:

\textbf{Experimento 1 (controlado):} 50 datasets de classificação --- 4 reais da biblioteca \texttt{scikit-learn} (Iris, Wine, Breast Cancer, Digits) e 46 sintéticos gerados com \texttt{make\_classification}, variando $N \in [200, 1000]$ instâncias, $D \in [6, 25]$ atributos, $C \in \{2, 3\}$ classes e Imbalance Ratio $\in [1{,}0, 20{,}0]$.

\textbf{Experimento 2 (realista inicial):} 144 datasets reais do OpenML~\cite{vanschoren2014openml}, selecionados por critérios de qualidade (sem valores ausentes, status ativo, dataset de referência).

\textbf{Experimento 3 (realista expandido):} 200 datasets reais do OpenML, com maior diversidade de tarefas de classificação. O perfil das meta-features é apresentado na Tabela~\ref{tab:metadataset_profile}.

\textbf{Experimento 4 (PyMFE):} os mesmos 200 datasets do Experimento~3, mas substituindo as 8 meta-features manuais por 78 descritores extraídos automaticamente pelo \texttt{PyMFE}~\cite{pymfe2020} (grupos \texttt{general}, \texttt{statistical}, \texttt{info-theory} e \texttt{clustering}), com imputação por mediana e recorte de valores superiores a $10^{10}$.

\begin{table}[ht]
\centering
\caption{Perfil das meta-features dos Experimentos 3 e 4 (200 datasets reais do OpenML).}
\label{tab:metadataset_profile}
\small
\begin{tabular}{lrrrrr}
\toprule
\textbf{Meta-feature} & \textbf{Mín.} & \textbf{Q1} & \textbf{Mediana} & \textbf{Q3} & \textbf{Máx.} \\
\midrule
\# Instâncias ($N$)   & 21    & 187    & 351,0   & 712,0    & 2.000 \\
\# Atributos ($D$)    & 2     & 9      & 18,0    & 37,0     & 2.836 \\
Razão $D/N$          & 0,002 & 0,023  & 0,048   & 0,096    & 7,144 \\
Imbalance Ratio (IR)  & 1,00  & 1,17   & 2,10    & 5,31     & 534,0 \\
Entropia              & 0,00  & 0,83   & 0,99    & 1,00     & 3,32  \\
Pearson médio         & 0,00  & 0,08   & 0,17    & 0,34     & 0,78  \\
\bottomrule
\end{tabular}
\end{table}

\subsection{Meta-features Extraídas}

Nos Experimentos 1--3, foram calculadas 8 meta-features estáticas: número de instâncias, número de atributos, razão $D/N$, Imbalance Ratio, entropia de Shannon das classes, correlação média de Pearson (valor absoluto) entre atributos, variância média dos atributos e informação mútua média dos atributos com o alvo.

No Experimento~4, essas 8 features foram substituídas pelos 78 descritores do \texttt{PyMFE}, cobrindo quatro grupos: \texttt{general} (dimensionalidade, proporção de classes), \texttt{statistical} (correlação canônica, assimetria, variância), \texttt{info-theory} (entropia, informação mútua, concentração) e \texttt{clustering} (Calinski-Harabasz, Davies-Bouldin, Dunn).

Em todos os experimentos, as meta-features de entrada $X$ foram padronizadas com StandardScaler ajustado exclusivamente no fold de treino.

\subsection{Algoritmos Base e Meta-alvo}

Cinco algoritmos base foram avaliados em cada dataset via validação cruzada estratificada de 5 folds, com F1-macro como métrica: \textbf{KNN} ($k=5$), \textbf{Decision Tree} (CART), \textbf{Random Forest} (100 estimadores), \textbf{Gaussian Naive Bayes} e \textbf{SVM} (kernel RBF). Todos os hiperparâmetros mantiveram os valores padrão do \texttt{scikit-learn} 1.2.2, com \textit{random state} fixado em 42.

\subsection{Meta-modelos}

Quatro meta-modelos foram avaliados como regressores multi-saída (prevendo o F1 dos 5 algoritmos simultaneamente):
\begin{itemize}
  \item \textbf{Ridge} ($\alpha = 1{,}0$): modelo linear;
  \item \textbf{SVR} (kernel RBF, $C=1{,}0$, $\varepsilon=0{,}1$): modelo não-linear com kernel implícito;
  \item \textbf{Random Forest Regressor} (100 estimadores): ensemble baseado em árvores;
  \item \textbf{MLP Regressor} (camadas $[16, 8]$, solver Adam, \texttt{max\_iter=2000}): rede neural.
\end{itemize}

\subsection{Protocolo de Validação}

Utilizou-se \textbf{validação cruzada de 5 folds no nível meta} (divisão por datasets). Em cada fold, o scaler do meta-alvo foi ajustado nos dados de treino, o meta-modelo foi treinado no alvo escalonado e as predições de teste foram invertidas para a escala original antes do cálculo das métricas --- garantindo que não ocorra vazamento de dados (\textit{data leakage}).

\subsection{Testes Estatísticos}

Para verificar a significância das diferenças entre scalers, empregamos:
\begin{itemize}
  \item \textbf{Teste de Friedman} (global): verifica se ao menos um scaler difere dos demais por meta-modelo, usando os vetores de MAE por dataset ($\alpha = 0{,}05$);
  \item \textbf{Wilcoxon signed-rank} (post-hoc): compara cada scaler contra o Baseline par a par.
\end{itemize}

% ============================================================
\section{Resultados e Discussão}
% ============================================================

\subsection{Distribuição do Meta-Alvo}

A Tabela~\ref{tab:base_perf} mostra os dados descritivos do meta-alvo nos 200 datasets reais. O desvio padrão alto indica que a variabilidade é grande: existem conjuntos de dados fáceis, onde o F1-score fica perto de 1,0, e outros bem difíceis, com F1-score perto de 0,0. Essa variação é o que torna o escalonamento tão importante.

\begin{table}[H]
\centering
\caption{Distribuição do F1-macro (meta-alvo) para os 5 algoritmos base nos 200 datasets reais.}
\label{tab:base_perf}
\small
\begin{tabular}{lcccc}
\toprule
\textbf{Algoritmo} & \textbf{Média} & \textbf{Desv.\ Pad.} & \textbf{Mediana} & \textbf{Máx.} \\
\midrule
KNN           & 0,654 & 0,197 & 0,643 & 1,000 \\
Decision Tree & 0,720 & 0,185 & 0,702 & 1,000 \\
Random Forest & 0,734 & 0,191 & 0,735 & 1,000 \\
Gaussian NB   & 0,643 & 0,210 & 0,648 & 1,000 \\
SVM           & 0,589 & 0,239 & 0,555 & 1,000 \\
\bottomrule
\end{tabular}
\end{table}

\subsection{Resultados por Experimento}

Para entender melhor as descobertas sem se perder em detalhes, reunimos as discussões dos quatro experimentos a seguir. Assim, conseguimos ver como o comportamento dos modelos mudou ao passarmos de dados sintéticos controlados para dados reais e complexos.

O Experimento 1 (com 50 datasets controlados) serviu como ponto de partida e confirmou algumas suposições: os modelos Ridge e Random Forest não sofreram impacto com escalonadores lineares, o SVR melhorou um pouco em termos de erro e ordenação, e o MLP falhou completamente quando os dados não foram centralizados.

Os Experimentos 2 e 3 (com 144 e 200 datasets reais e 8 meta-features) confirmaram esses comportamentos em cenários de dados reais e variados. O colapso do MLP foi ainda mais nítido aqui: no Experimento 2, a correlação de Spearman despencou para 0,034 (indicando nenhuma capacidade de ordenação) e a acurácia de recomendação Acc@1 caiu para 12,5\% (o que é pior que o acaso de 20\% para 5 opções). Em contrapartida, usar StandardScaler ou RobustScaler recuperou a capacidade de previsão do modelo, com o MAE caindo para a faixa de 0,11 a 0,13 e a correlação de Spearman subindo para valores entre 0,41 e 0,45. A Tabela~\ref{tab:mlp_comparison} na Seção~\ref{subsec:analise_mlp} resume o desempenho do MLP ao longo dos quatro testes.

O Experimento 4 (com 200 datasets e 78 meta-features do PyMFE) traz os resultados do pipeline mais completo do estudo. O uso dessas características mais ricas reduziu bastante o erro de previsão de todos os modelos. Por exemplo, o menor MAE obtido pelo SVR caiu de 0,123 (no Experimento 3) para 0,088 (no Experimento 4), e o do Random Forest baixou de 0,114 para 0,085. A Tabela~\ref{tab:results_exp4_summary} resume os resultados desse cenário final, comparando as baselines com as melhores configurações de escalonamento.

\begin{table}[ht]
\centering
\caption{Resumo dos resultados do Experimento 4 (PyMFE, 200 datasets). Comparativo entre Baseline e escalonadores recomendados.}
\label{tab:results_exp4_summary}
\small
\begin{tabular}{llccc}
\toprule
\textbf{Meta-Modelo} & \textbf{Técnica} & \textbf{MAE} & \textbf{Spearman} & \textbf{Acc@1} \\
\midrule
\multirow{2}{*}{Ridge}
 & Baseline / Scalers Lineares & \textbf{0,1007} & \textbf{0,4966} & \textbf{0,3788} \\
 & QuantileTransformer         & 0,1152          & 0,4483          & 0,3586          \\
\midrule
\multirow{3}{*}{SVR}
 & Baseline                    & 0,0978          & 0,5213          & 0,3687          \\
 & StandardScaler (Melhor Rank)& 0,0898          & \textbf{0,5982} & \textbf{0,4949} \\
 & RobustScaler (Melhor MAE)   & \textbf{0,0884} & 0,5966          & 0,4646          \\
\midrule
\multirow{2}{*}{Random Forest}
 & Baseline                    & 0,0860          & 0,6019          & 0,4747          \\
 & MinMaxScaler (Melhor Geral) & \textbf{0,0845} & \textbf{0,6157} & \textbf{0,5152} \\
\midrule
\multirow{4}{*}{MLP}
 & Baseline                    & 0,1830          & 0,0802          & 0,2121          \\
 & StandardScaler (Melhor MAE) & \textbf{0,1165} & 0,3875          & 0,3384          \\
 & RobustScaler (Melhor Acc@1) & 0,1220          & 0,3542          & \textbf{0,3586} \\
 & QuantileTransformer (Melhor $\rho_s$) & 0,1294 & \textbf{0,4355} & 0,2980          \\
\bottomrule
\end{tabular}
\end{table}

\subsection{Testes de Hipóteses}

\textbf{Teste de Friedman.} Este teste analisa se há alguma diferença estatisticamente significativa no erro absoluto (MAE) dos modelos ao variarmos os escalonadores (os resultados estão na Tabela~\ref{tab:friedman_all}). Nos Experimentos 2 e 3, apenas o MLP mostrou diferenças significativas, o que indica que os modelos Ridge, SVR e Random Forest conseguem lidar bem com os dados brutos. No entanto, no Experimento 4 (com PyMFE), a maior quantidade de características deu mais precisão para a análise: o Ridge e o SVR também passaram a apresentar diferenças estatísticas significativas. Por outro lado, o Random Forest continuou muito estável e não sofreu impacto estatístico sob nenhuma técnica.

\begin{table}[H]
\centering
\caption{Resultados do teste de Friedman por meta-modelo nos Experimentos~2, 3 e 4.}
\label{tab:friedman_all}
\small
\begin{tabular}{llccl}
\toprule
\textbf{Experimento} & \textbf{Meta-Modelo} & \textbf{$\chi^2$} & \textbf{$p$-value} & \textbf{Decisão} \\
\midrule
\multirow{4}{*}{Exp.~2 (144 ds)}
 & Ridge        & 3,04            & 0,551                                    & Não significativo \\
 & SVR          & 8,37            & 0,079                                    & Não significativo \\
 & Random Forest & 5,41           & 0,248                                    & Não significativo \\
 & MLP          & \textbf{148,72} & $\mathbf{3{,}8 \times 10^{-31}}$        & \textbf{Significativo} \\
\midrule
\multirow{4}{*}{Exp.~3 (200 ds)}
 & Ridge        & 1,61            & 0,807                                    & Não significativo \\
 & SVR          & 3,41            & 0,491                                    & Não significativo \\
 & Random Forest & 0,34           & 0,987                                    & Não significativo \\
 & MLP          & \textbf{110,71} & $\mathbf{5{,}1 \times 10^{-23}}$        & \textbf{Significativo} \\
\midrule
\multirow{4}{*}{Exp.~4 (PyMFE)}
 & Ridge        & \textbf{22,38}  & $\mathbf{1{,}7 \times 10^{-4}}$         & \textbf{Significativo} \\
 & SVR          & \textbf{53,92}  & $\mathbf{5{,}5 \times 10^{-11}}$        & \textbf{Significativo} \\
 & Random Forest & 4,58           & 0,333                                    & Não significativo \\
 & MLP          & \textbf{278,02} & $\mathbf{5{,}9 \times 10^{-59}}$        & \textbf{Significativo} \\
\bottomrule
\end{tabular}
\end{table}

\textbf{Wilcoxon post-hoc para o MLP.} Os testes de comparação par a par (Wilcoxon) confirmaram que qualquer escalonamento é muito melhor do que deixar o meta-alvo do MLP sem tratamento nos três experimentos com dados reais ($p < 10^{-4}$). O StandardScaler e o RobustScaler diminuíram o erro absoluto em até 36\% e ajudaram o modelo a ordenar melhor as opções (a correlação de Spearman subiu de valores quase nulos para mais de 0,40). Por exemplo, no Experimento 4, o StandardScaler reduziu o MAE do MLP de 0,183 para 0,117 e elevou a correlação de Spearman de 0,080 para 0,388.

\subsection{Análise por Meta-Modelo}

\subsubsection{Ridge Regression: Invariância a Transformações Lineares}

O modelo Ridge apresentou resultados \textit{idênticos} para os escalonadores lineares (StandardScaler, MinMaxScaler e RobustScaler) em todos os quatro experimentos. Isso acontece por conta de uma característica matemática simples: quando aplicamos uma transformação linear como $\tilde{y} = a \cdot y + b$, o Ridge apenas ajusta os coeficientes para compensar a escala. Ao fazer a inversão do escalonamento na saída, o resultado final volta a ser idêntico ao original. A única exceção é o QuantileTransformer. Como ele aplica uma transformação não-linear, a ordem dos dados no espaço-alvo é distorcida e a inversão não consegue recuperar a escala original, o que explica a queda na acurácia de recomendação (Acc@1) observada em todos os cenários.

\subsubsection{SVR: O Escalonamento Ajuda Mais na Recomendação do que na Redução do Erro}

Para o SVR, o escalonamento faz uma grande diferença na qualidade das recomendações (Acc@1 e correlação de Spearman) em todos os cenários, embora o impacto no erro absoluto (MAE) varie de acordo com o experimento. Em dados reais sem o PyMFE (Experimentos 2 e 3), o erro mudou muito pouco, mas a capacidade de recomendar o melhor algoritmo aumentou significativamente: no Experimento 2, o Acc@1 subiu de 0,486 para 0,563 com RobustScaler, e no Experimento 3 foi de 0,429 para 0,470. Já no Experimento 4, com as meta-features mais completas do PyMFE, o modelo melhorou em todas as frentes: o MAE caiu cerca de 10\% (de 0,098 para 0,088 com RobustScaler) e a acurácia de recomendação foi de 0,369 para 0,495 com StandardScaler. 

\subsubsection{Random Forest: Robustez ao Escalonamento}

O Random Forest confirmou ser o modelo mais resistente a mudanças de escala no meta-alvo. Em todos os experimentos, a variação de MAE entre as técnicas lineares foi minúscula (menor que 0,007), e os testes estatísticos não indicaram diferença significativa global. Essa estabilidade acontece pela própria natureza dos modelos de árvore, que tomam decisões de divisão baseadas na ordenação relativa (ranks) dos dados, mantida por qualquer transformação linear. A única técnica que atrapalha o Random Forest é o QuantileTransformer: por alterar a ordem relativa dos dados, ele prejudicou a capacidade de recomendação (Acc@1) em todos os cenários testados.

\subsubsection{MLP: Colapso sem Escalonamento Centralizado}
\label{subsec:analise_mlp}

O resultado mais marcante é o comportamento do MLP. Sem usar um escalonador centralizado no meta-alvo, a rede neural simplesmente falha em aprender qualquer padrão útil nos dados reais. A correlação de Spearman cai para patamares quase nulos (0,034 no Exp. 2, 0,167 no Exp. 3 e 0,080 no Exp. 4) e a acurácia de recomendação fica em torno ou abaixo do acaso de 20\% (chegando a apenas 12,5\% no Exp. 2). A Tabela~\ref{tab:mlp_comparison} detalha esse colapso nos quatro cenários:

\begin{table}[ht]
\centering
\caption{Comportamento do MLP nos quatro experimentos: Baseline versus melhor scaler.}
\label{tab:mlp_comparison}
\small
\begin{tabular}{lcccc}
\toprule
\textbf{Experimento} & \textbf{Scaler} & \textbf{MAE} & \textbf{Spearman} & \textbf{Acc@1} \\
\midrule
\multirow{2}{*}{Exp.\ 1 (sintético, 50 ds)} & Baseline     & 0,147 & 0,274 & 0,320 \\
                                             & RobustScaler & \textbf{0,071} & \textbf{0,588} & \textbf{0,440} \\
\midrule
\multirow{2}{*}{Exp.\ 2 (real, 144 ds)}     & Baseline     & 0,179 & 0,034 & 0,125 \\
                                             & RobustScaler & \textbf{0,118} & \textbf{0,453} & \textbf{0,438} \\
\midrule
\multirow{3}{*}{Exp.\ 3 (real, 200 ds)}     & Baseline     & 0,174 & 0,167 & 0,207 \\
                                             & StandardScaler & \textbf{0,131} & \textbf{0,431} & 0,374 \\
                                             & RobustScaler & 0,135 & 0,413 & \textbf{0,399} \\
\midrule
\multirow{3}{*}{Exp.\ 4 (PyMFE, 200 ds)}    & Baseline     & 0,183 & 0,080 & 0,212 \\
                                             & QuantileTransformer & \textbf{0,129} & \textbf{0,436} & 0,298 \\
                                             & RobustScaler & 0,122 & 0,354 & \textbf{0,359} \\
\bottomrule
\end{tabular}
\end{table}

Esse colapso é muito forte com dados reais porque o meta-alvo apresenta alta variância e alguns valores extremos de desempenho (algoritmos com F1-score próximo de 0 ou de 1 no mesmo conjunto de dados). Quando esses dados não são centralizados em torno de zero, os pesos da rede neural sofrem com gradientes instáveis ou saturação das funções de ativação. Técnicas como StandardScaler e RobustScaler salvam o modelo porque centralizam e limitam a dispersão dos dados.

O MinMaxScaler não ajuda o MLP porque ele apenas mapeia os dados para o intervalo $[0, 1]$, mantendo a assimetria e a falta de centralização (o que explica por que ele dá resultados quase idênticos ao Baseline sem escalonamento). Por fim, no Experimento 4 (com as 78 características do PyMFE), um comportamento curioso: o QuantileTransformer deu a melhor correlação de Spearman do MLP (0,436). Isso mostra que, em espaços com características mais densas e estruturadas, essa transformação não-linear por quantis pode ajudar a rede neural a identificar melhor a hierarquia entre os algoritmos base. 

\subsection{O Paradoxo do QuantileTransformer}

Embora o QuantileTransformer possa ajudar a reduzir o erro absoluto (MAE) ao transformar os dados em uma distribuição normal, ele costuma piorar a capacidade de recomendação (Acc@1) dos modelos. Isso acontece porque a sua transformação não-linear muda as distâncias relativas entre os scores de F1 dos algoritmos. Por exemplo, se dois algoritmos têm desempenhos muito parecidos (como 0,70 e 0,71), o QuantileTransformer pode afastar bastante os seus valores transformados, enquanto dois desempenhos distantes (como 0,50 e 0,70) podem ficar mais próximos. Essa distorção dificulta que a inversão do escalonamento recupere a ordem exata das previsões.

A lição prática é: se a intenção do sistema é de fato \textbf{recomendar o melhor algoritmo}, o QuantileTransformer deve ser evitado para a maioria dos modelos. A única exceção notável foi o MLP no Experimento 4 (com PyMFE), no qual o uso de características mais ricas na entrada parece ter ajudado a contornar esse problema.

\subsection{Consistência entre os Experimentos}

A Tabela~\ref{tab:consistency} sintetiza a consistência dos achados ao longo dos quatro experimentos, mostrando que os padrões identificados com dados controlados são robustos a cenários reais heterogêneos.

\begin{table}[ht]
\centering
\caption{Consistência dos achados ao longo dos quatro experimentos.}
\label{tab:consistency}
\resizebox{\columnwidth}{!}{%
\begin{tabular}{lccccc}
\toprule
\textbf{Fenômeno observado} & \textbf{Exp.\ 1} & \textbf{Exp.\ 2} & \textbf{Exp.\ 3} & \textbf{Exp.\ 4} & \textbf{Avaliação} \\
\midrule
Ridge invariante a scalers lineares        & \checkmark & \checkmark & \checkmark & \checkmark & Alta \\
Ridge degradado pelo QT (Acc@1)            & \checkmark & \checkmark & \checkmark & \checkmark & Alta \\
MLP: colapso sem escalonamento centralizado & Parcial    & Total      & Total      & Total      & Reforçada \\
MLP: ganho no MAE com SS/RS                & 51,8\%     & 34,3\%     & 24,9\%     & 36,3\%     & Alta \\
MLP: MinMaxScaler sem ganho                & \checkmark & \checkmark & \checkmark & \checkmark & Alta \\
RF invariante a scalers lineares            & \checkmark & \checkmark & \checkmark & \checkmark & Alta \\
QT degrada rankings em todos os modelos     & \checkmark & \checkmark & \checkmark & \checkmark & Alta \\
SVR beneficiado pelo escalonamento         & MAE        & Spearman   & Spearman   & Todos      & Alta \\
\bottomrule
\end{tabular}
}%
\end{table}

\section{Conclusão}
% ============================================================

Este estudo analisou como o escalonamento do meta-alvo influencia o desempenho de sistemas de meta-aprendizagem na tarefa de selecionar algoritmos. Para isso, realizaram-se quatro experimentos: 50 datasets em ambiente controlado, 144 datasets reais (etapa inicial), 200 datasets reais (etapa expandida) e, finalmente, 200 datasets reais utilizando 78 meta-features extraídas automaticamente pelo PyMFE.

Os resultados deixam claro que \textbf{a sensibilidade ao escalonamento do meta-alvo depende muito mais da arquitetura do meta-modelo} do que do tipo ou da origem dos dados analisados:

\begin{itemize}
  \item \textbf{Ridge}: é imune a transformações lineares no alvo por conta de sua formulação matemática. Apenas escalonamentos não-lineares conseguem alterar o seu comportamento.
  \item \textbf{SVR}: melhora bastante a qualidade das recomendações (acurácia Acc@1 e Spearman) ao usar StandardScaler ou RobustScaler. Essa vantagem fica ainda mais nítida e estatisticamente significativa à medida que se usa características de entrada mais ricas (como as do PyMFE).
  \item \textbf{Random Forest}: mostrou-se muito estável e quase não sofreu impactos com os escalonadores lineares, sendo o modelo mais seguro para uso geral. Curiosamente, ao combinar o Random Forest com MinMaxScaler no Experimento 4, obteve-se a maior correlação de Spearman de todo o estudo (0,616).
  \item \textbf{MLP}: o escalonamento centralizado (usando StandardScaler ou RobustScaler) é obrigatório. Sem esse ajuste, a rede neural falha nos cenários reais, dando palpites equivalentes a um sorteio aleatório (a correlação de Spearman fica abaixo de 0,09 no Experimento 4).
\end{itemize}

A principal recomendação geral é que o \textbf{QuantileTransformer costuma prejudicar a qualidade das recomendações} (Acc@1) na maioria dos meta-modelos, sendo a única técnica não recomendada para seleção de algoritmos — embora ele tenha melhorado a Spearman do MLP no Experimento 4 com PyMFE.

A adoção de meta-features mais completas (PyMFE) reduziu drasticamente o erro de previsão em todos os modelos.

Por fim, confirma-se a hipótese H3 de que não existe um escalonador perfeito para todas as situações. A melhor escolha sempre vai depender do modelo que você está usando e do seu objetivo principal (seja reduzir o erro numérico ou obter as melhores recomendações). Como regra geral e prática, recomenda-se o uso do \textbf{RobustScaler}, pois ele oferece o melhor equilíbrio entre precisão e estabilidade, especialmente ao lidar com dados reais e que possuem desempenhos muito variados.

\bibliographystyle{sbc}
\bibliography{relatorio-sbc}

\end{document}
